\subsubsection{Manual}\label{trien-khai-remediate-manual}

Khắc phục thủ công được áp dụng đối với các khuyến nghị về thiết kế hạ tầng hoặc yêu cầu can thiệp sâu của quản trị viên hệ thống để tránh làm gián đoạn dịch vụ.

\textbf{Bảng chi tiết các hoạt động Remediation thủ công:}

{\def\LTcaptype{none} % do not increment counter
\begin{longtable}[]{@{}|
  >{\raggedright\arraybackslash}p{(\linewidth - 8\tabcolsep - 5\arrayrulewidth) * \real{0.0739}}|
  >{\raggedright\arraybackslash}p{(\linewidth - 8\tabcolsep - 5\arrayrulewidth) * \real{0.2995}}|
  >{\raggedright\arraybackslash}p{(\linewidth - 8\tabcolsep - 5\arrayrulewidth) * \real{0.3000}}|
  >{\raggedright\arraybackslash}p{(\linewidth - 8\tabcolsep - 5\arrayrulewidth) * \real{0.3266}}|@{}}
\hline
\begin{minipage}[b]{\linewidth}\raggedright
\textbf{Phần}
\end{minipage} & \begin{minipage}[b]{\linewidth}\raggedright
\textbf{Benchmark}
\end{minipage} & \begin{minipage}[b]{\linewidth}\raggedright
\textbf{Phương pháp thực hiện}
\end{minipage} & \begin{minipage}[b]{\linewidth}\raggedright
\textbf{Hành động thực hiện}
\end{minipage} \\
\hline
\endhead

\multirow{19}{=}{1} & Phân vùng riêng cho các container & Tạo phân vùng LVM riêng và định dạng ext4 & Dừng Docker, backup dữ liệu, mount phân vùng mới, cập nhật /etc/fstab, và khởi động lại Docker \\
\cline{2-4}
& Chỉ cho trusted user kiểm soát Docker Daemon & Sử dụng gpasswd để xóa các user không được phép khỏi nhóm docker & Chạy lệnh sudo gpasswd -d <username> docker để gỡ bỏ user \\
\cline{2-4}
& Cấu hình kiểm tra nhật ký cho Docker Daemon & Cấu hình audit rule cho dockerd trong docker.rules & Ghi cấu hình -w /usr/bin/dockerd -k docker và reload auditd \\
\cline{2-4}
& Audit folder /run/containerd & Cấu hình audit rule cho /run/containerd trong docker.rules & Ghi cấu hình -a exit,always -F path=/run/containerd -F perm=war -k docker và reload auditd \\
\cline{2-4}
& Audit folder /var/lib/docker & Cấu hình audit rule cho /var/lib/docker trong docker.rules & Ghi cấu hình -a exit,always -F path=/var/lib/docker -F perm=war -k docker và reload auditd \\
\cline{2-4}
& Audit folder /etc/Docker & Cấu hình audit rule cho /etc/docker trong docker.rules & Ghi cấu hình -w /etc/docker -k docker và reload auditd \\
\cline{2-4}
& Audit docker.service & Cấu hình audit rule cho docker.service trong docker.rules & \begin{minipage}[t]{\linewidth}\raggedright
Ghi cấu hình -w /usr/lib/systemd/\\
system/docker.service -k docker\strut
\end{minipage} \\
\cline{2-4}
& Audt containerd.sock & Cấu hình audit rule cho containerd.sock trong docker.rules & Ghi cấu hình -w /run/containerd/containerd.sock -k docker \\
\cline{2-4}
& Audit docker.sock & Cấu hình audit rule cho docker.sock trong docker.rules & Ghi cấu hình -w /run/docker.sock -k docker \\
\cline{2-4}
& Audt /etc/default/docker & Cấu hình audit rule cho /etc/default/docker trong docker.rules & Ghi cấu hình -w /etc/default/docker -k docker \\
\cline{2-4}
& Audit /etc/docker/daemon.json & Cấu hình audit rule cho /etc/docker/daemon.json trong docker.rules & Ghi cấu hình -w /etc/docker/daemon.json -k docker \\
\cline{2-4}
& Audit /etc/containerd/config.toml & Cấu hình audit rule cho /etc/containerd/config.toml & Ghi cấu hình -w /etc/containerd/config.toml -k docker \\
\cline{2-4}
& Aduti /etc/sysconfig/docker & Cấu hình audit rule cho /etc/sysconfig/docker & Ghi cấu hình -w /etc/sysconfig/docker -k docker \\
\cline{2-4}
& Audit /usr/bin/containerd-shim & Cấu hình audit rule cho /usr/bin/containerd-shim & Ghi cấu hình -w /usr/bin/containerd-shim -k docker \\
\cline{2-4}
& Audit /usr/bin/containerd-shim-runc-v1 & Cấu hình audit rule cho containerd-shim-runc-v1 & Ghi cấu hình -w /usr/bin/containerd-shim-runc-v1 -k docker \\
\cline{2-4}
& Audit /usr/bin/containerd-shim-runc-v2 & Cấu hình audit rule cho containerd-shim-runc-v2 & Ghi cấu hình -w /usr/bin/containerd-shim-runc-v2 -k docker \\
\cline{2-4}
& Audit /usr/bin/runc & Cấu hình audit rule cho runc & Ghi cấu hình -w /usr/bin/runc -k docker \\
\cline{2-4}
& Đảm báo bảo mật cho container & Cấu hình hardening OS của container host & Cập nhật OS, tắt dịch vụ thừa, cấu hình SSH và tường lửa \\
\cline{2-4}
& Đảm bảo phiên bản mới cho Docker kt & Cập nhật Docker lên phiên bản mới nhất & Cập nhật danh sách gói và chạy apt install --only-upgrade docker-ce \\
\hline
\multirow{19}{=}{2} & Chạy docker daemon tránh quyền root & & \\
\cline{2-4}
& Network trafic được đặt hạn chế giữa các container ở default bridge & & \\
\cline{2-4}
& Logging level set ở ``info'' & & \\
\cline{2-4}
& Docker được quyền thay đổi iptables & & \\
\cline{2-4}
& Đảm bảo không sử dụng registry không an toàn & & \\
\cline{2-4}
& Đảm bảo không sử dụng trình điều khiển aufs & & \\
\cline{2-4}
& Đảm bảo không dùng devicemapper & & \\
\cline{2-4}
& Đảm bảo xác thực TLS cho Docker Daemon & & \\
\cline{2-4}
& Đảm bảo giới hạn ulimit mặc định & & \\
\cline{2-4}
& Hỗ trợ user namespace & & \\
\cline{2-4}
& Cgroup sử dụng được xác nhận & & \\
\cline{2-4}
& Đảm bảo device size không thay đổi trừ khi cần thiết & & \\
\cline{2-4}
& Đảm bảo authorization cho Docker client & & \\
\cline{2-4}
& Đảm bảo container bj hạn chế không cho quyền mới & & \\
\cline{2-4}
& Đảm bảo container có live restore & & \\
\cline{2-4}
& Đảm bảo có log cho đăng nhập gần và từ xa & & \\
\cline{2-4}
& Đảm bảo Userland Proxy đã tắt & & \\
\cline{2-4}
& Đảm bảo daemon-wide custom seccomp profile được thiết đặt & & \\
\cline{2-4}
& Đảm bảo các tính năng thử nghiệm không được triển khai trong môi trường sản xuất & & \\
\hline
\multirow{14}{=}{3} & Đảm bảo quyền root:root cho docker.service & & \\
\cline{2-4}
& Đảm bảo docker.service cho quyền hợp lý & & \\
\cline{2-4}
& Đảm bảo quyền sở hữu của docker.socket được đặt thành root:root & & \\
\cline{2-4}
& Đảm bảo quyền truy cập tệp docker.socket được đặt thành 644 hoặc hạn chế & & \\
\cline{2-4}
& Đảm bảo quyền sỡ hữu thư mục /etc/docker được đặt thành root:root & & \\
\cline{2-4}
& Đảm bảo quyền truy cập thư mục /etc/docker được đặt thành 755 hoặc hạn chế & & \\
\cline{2-4}
& Đảm bảo quyền truy cập tệ docker.socker được đặt thành 644 hoặc hạn chế & & \\
\cline{2-4}
& Đảm bảo quyền root:root cho sỡ hữu tệp chứng chỉ & & \\
\cline{2-4}
& Đảm bảo quyền truy cập chứng trí registry được đặt thành 444 hoặc hạn chế hơn & & \\
\cline{2-4}
& Đảm bảo quyền sỡ hữu Docker server là root:root và file permission 4444 & & \\
\cline{2-4}
& Đảm bảo quyền Docker socker file ownership là root:docker và để file permission 660 & & \\
\cline{2-4}
& Đảm bảo quyền daemon.json là root:root và file permssion 644 & & \\
\cline{2-4}
& Đảm bảo quyền /etc/sysconfig/docker là root:root và file permission 644 & & \\
\cline{2-4}
& Đảm bảo Containered socket có quyền là root:root và file permission 660 & & \\
\hline
\multirow{9}{=}{4} & Phải có container tạo, chạy trong trusted base images và không có package dư thừa không cần thiết được tải về. & & \\
\cline{2-4}
& Đảm bảo image được scan và phải có security patches & & \\
\cline{2-4}
& Phải có HEALTHCHECK instruction & & \\
\cline{2-4}
& Phải có update instruction và không sử dụng một mình trong Dockerfiles & & \\
\cline{2-4}
& Không cung cấp quyền Setuid và Getuid & & \\
\cline{2-4}
& Dùng COPY thay vì ADD & & \\
\cline{2-4}
& Secrets không được bỏ vào Dockerfiles & & \\
\cline{2-4}
& Các package phải chính thống khi tải về & & \\
\cline{2-4}
& Các signed artifacts phải được kiểm chứng & & \\
\hline
\multirow{32}{=}{5} & 5.1: Khong bat swarm mode neu khong can thiet & \begin{minipage}[t]{\linewidth}\raggedright
Roi Swarm neu host khong dung cluster.\strut
\end{minipage} & \begin{minipage}[t]{\linewidth}\raggedright
Chay docker swarm leave --force sau khi xac nhan khong con workload Swarm.\strut
\end{minipage} \\
\cline{2-4}
 & 5.2: Bat AppArmor Profile neu moi truong ho tro & \begin{minipage}[t]{\linewidth}\raggedright
Khoi chay container voi AppArmor profile phu hop.\strut
\end{minipage} & \begin{minipage}[t]{\linewidth}\raggedright
Dung security-opt apparmor profile hoac profile mac dinh da phe duyet.\strut
\end{minipage} \\
\cline{2-4}
 & 5.3: Thiet lap SELinux security options neu ap dung & \begin{minipage}[t]{\linewidth}\raggedright
Chi ap dung SELinux label khi host thuc su bat SELinux.\strut
\end{minipage} & \begin{minipage}[t]{\linewidth}\raggedright
Bat SELinux dung moi truong hoac bo label khong co tac dung tren Ubuntu/AppArmor.\strut
\end{minipage} \\
\cline{2-4}
 & 5.4: Han che Linux kernel capabilities & \begin{minipage}[t]{\linewidth}\raggedright
Drop capability mac dinh va chi add lai capability can thiet.\strut
\end{minipage} & \begin{minipage}[t]{\linewidth}\raggedright
Dung cap-drop ALL; kiem thu roi chi cap-add capability da duoc phe duyet.\strut
\end{minipage} \\
\cline{2-4}
 & 5.5: Khong su dung privileged container & \begin{minipage}[t]{\linewidth}\raggedright
Tai tao container khong dung privileged mode.\strut
\end{minipage} & \begin{minipage}[t]{\linewidth}\raggedright
Bo privileged va thay bang capability hoac device cu the neu bat buoc.\strut
\end{minipage} \\
\cline{2-4}
 & 5.6: Khong mount thu muc he thong nhay cam cua host & \begin{minipage}[t]{\linewidth}\raggedright
Loai bo bind mount nhay cam.\strut
\end{minipage} & \begin{minipage}[t]{\linewidth}\raggedright
Bo mount /, /etc, /proc, /sys, /dev; neu bat buoc thi read-only va co phe duyet.\strut
\end{minipage} \\
\cline{2-4}
 & 5.7: Khong chay sshd trong container & \begin{minipage}[t]{\linewidth}\raggedright
Loai bo SSH daemon khoi image/container.\strut
\end{minipage} & \begin{minipage}[t]{\linewidth}\raggedright
Quan tri bang docker exec, logs hoac orchestration thay vi sshd trong container.\strut
\end{minipage} \\
\cline{2-4}
 & 5.8: Khong anh xa privileged port khong can thiet & \begin{minipage}[t]{\linewidth}\raggedright
Go privileged port hoac ghi nhan phe duyet.\strut
\end{minipage} & \begin{minipage}[t]{\linewidth}\raggedright
Dung port cao hon hoac reverse proxy; chi giu port nho hon 1024 khi co nhu cau ro rang.\strut
\end{minipage} \\
\cline{2-4}
 & 5.9: Chi mo port can thiet & \begin{minipage}[t]{\linewidth}\raggedright
Xoa cac port publish khong dung.\strut
\end{minipage} & \begin{minipage}[t]{\linewidth}\raggedright
Cap nhat Docker CLI/Compose/Ansible de chi publish port nam trong thiet ke.\strut
\end{minipage} \\
\cline{2-4}
 & 5.10: Khong chia se network namespace cua host & \begin{minipage}[t]{\linewidth}\raggedright
Recreate container tren network rieng.\strut
\end{minipage} & \begin{minipage}[t]{\linewidth}\raggedright
Bo network host; dung bridge/overlay user-defined network.\strut
\end{minipage} \\
\cline{2-4}
 & 5.11: Gioi han memory cho container & \begin{minipage}[t]{\linewidth}\raggedright
Dat memory limit theo workload.\strut
\end{minipage} & \begin{minipage}[t]{\linewidth}\raggedright
Dung memory limit trong Docker CLI, playbook hoac compose.\strut
\end{minipage} \\
\cline{2-4}
 & 5.12: Dat CPU priority phu hop & \begin{minipage}[t]{\linewidth}\raggedright
Dat CPU shares hoac quota.\strut
\end{minipage} & \begin{minipage}[t]{\linewidth}\raggedright
Dung cpu-shares hoac CPU quota phu hop, sau do test hieu nang.\strut
\end{minipage} \\
\cline{2-4}
 & 5.13: Root filesystem o che do read-only & \begin{minipage}[t]{\linewidth}\raggedright
Chay container voi rootfs read-only.\strut
\end{minipage} & \begin{minipage}[t]{\linewidth}\raggedright
Dung read-only va khai bao tmpfs/volume rieng cho duong dan can ghi.\strut
\end{minipage} \\
\cline{2-4}
 & 5.14: Bind traffic vao interface cu the & \begin{minipage}[t]{\linewidth}\raggedright
Bind port vao IP duoc phe duyet.\strut
\end{minipage} & \begin{minipage}[t]{\linewidth}\raggedright
Dung 127.0.0.1:port:port hoac IP NIC cu the thay vi 0.0.0.0.\strut
\end{minipage} \\
\cline{2-4}
 & 5.15: Gioi han restart policy & \begin{minipage}[t]{\linewidth}\raggedright
Doi restart policy sang on-failure co gioi han.\strut
\end{minipage} & \begin{minipage}[t]{\linewidth}\raggedright
Dung restart on-failure:5 hoac restart policy/retry tuong ung.\strut
\end{minipage} \\
\cline{2-4}
 & 5.16: Khong chia se PID namespace cua host & \begin{minipage}[t]{\linewidth}\raggedright
Recreate container khong dung host PID.\strut
\end{minipage} & \begin{minipage}[t]{\linewidth}\raggedright
Bo pid host khoi cau hinh runtime.\strut
\end{minipage} \\
\cline{2-4}
 & 5.17: Khong chia se IPC namespace cua host & \begin{minipage}[t]{\linewidth}\raggedright
Recreate container khong dung host IPC.\strut
\end{minipage} & \begin{minipage}[t]{\linewidth}\raggedright
Bo ipc host; neu can chia se IPC thi gioi han giua container duoc phe duyet.\strut
\end{minipage} \\
\cline{2-4}
 & 5.18: Khong expose truc tiep host device & \begin{minipage}[t]{\linewidth}\raggedright
Go device mapping khong can thiet.\strut
\end{minipage} & \begin{minipage}[t]{\linewidth}\raggedright
Bo device mapping hoac gioi han dung thiet bi/quyen truy cap theo nghiep vu.\strut
\end{minipage} \\
\cline{2-4}
 & 5.19: Ghi de ulimit khi can & \begin{minipage}[t]{\linewidth}\raggedright
Khai bao ulimit ro rang.\strut
\end{minipage} & \begin{minipage}[t]{\linewidth}\raggedright
Dung ulimit hoac default-ulimits trong daemon.json da duoc harden.\strut
\end{minipage} \\
\cline{2-4}
 & 5.20: Khong dung shared mount propagation & \begin{minipage}[t]{\linewidth}\raggedright
Chuyen mount propagation sang private/slave.\strut
\end{minipage} & \begin{minipage}[t]{\linewidth}\raggedright
Cap nhat bind propagation va recreate container.\strut
\end{minipage} \\
\cline{2-4}
 & 5.21: Khong chia se UTS namespace cua host & \begin{minipage}[t]{\linewidth}\raggedright
Recreate container khong dung host UTS.\strut
\end{minipage} & \begin{minipage}[t]{\linewidth}\raggedright
Bo uts host khoi cau hinh runtime.\strut
\end{minipage} \\
\cline{2-4}
 & 5.22: Khong tat default seccomp profile & \begin{minipage}[t]{\linewidth}\raggedright
Bo cau hinh seccomp unconfined.\strut
\end{minipage} & \begin{minipage}[t]{\linewidth}\raggedright
Dung default seccomp hoac custom seccomp profile da kiem soat.\strut
\end{minipage} \\
\cline{2-4}
 & 5.23: Khong dung docker exec voi privileged option & \begin{minipage}[t]{\linewidth}\raggedright
Cam quy trinh van hanh dung exec privileged.\strut
\end{minipage} & \begin{minipage}[t]{\linewidth}\raggedright
Cap nhat runbook, review quyen admin va bat logging/audit tap trung.\strut
\end{minipage} \\
\cline{2-4}
 & 5.24: Khong dung docker exec voi user root option & \begin{minipage}[t]{\linewidth}\raggedright
Van hanh bang user ung dung.\strut
\end{minipage} & \begin{minipage}[t]{\linewidth}\raggedright
Dung user ung dung hoac sua image de co USER non-root.\strut
\end{minipage} \\
\cline{2-4}
 & 5.25: Xac nhan cgroup usage & \begin{minipage}[t]{\linewidth}\raggedright
Bo custom cgroup parent neu khong can.\strut
\end{minipage} & \begin{minipage}[t]{\linewidth}\raggedright
Loai bo cgroup parent hoac ghi tai lieu phe duyet kem resource control.\strut
\end{minipage} \\
\cline{2-4}
 & 5.26: Ngan container nhan them dac quyen & \begin{minipage}[t]{\linewidth}\raggedright
Them no-new-privileges.\strut
\end{minipage} & \begin{minipage}[t]{\linewidth}\raggedright
Dung security option no-new-privileges=true hoac security opts trong Ansible.\strut
\end{minipage} \\
\cline{2-4}
 & 5.27: Co runtime health check & \begin{minipage}[t]{\linewidth}\raggedright
Sua healthcheck de phan anh trang thai that.\strut
\end{minipage} & \begin{minipage}[t]{\linewidth}\raggedright
Them HEALTHCHECK hoac healthcheck command noi bo; nghiem thu khi healthy.\strut
\end{minipage} \\
\cline{2-4}
 & 5.28: Ghim phien ban image, khong dung latest & \begin{minipage}[t]{\linewidth}\raggedright
Ghim image bang version tag hoac digest.\strut
\end{minipage} & \begin{minipage}[t]{\linewidth}\raggedright
Doi nginx latest thanh version tag hoac sha256 digest da xac minh.\strut
\end{minipage} \\
\cline{2-4}
 & 5.29: Su dung PIDs cgroup limit & \begin{minipage}[t]{\linewidth}\raggedright
Dat PID limit.\strut
\end{minipage} & \begin{minipage}[t]{\linewidth}\raggedright
Dung pids limit trong Docker CLI, playbook hoac compose.\strut
\end{minipage} \\
\cline{2-4}
 & 5.30: Khong dung default bridge docker0 & \begin{minipage}[t]{\linewidth}\raggedright
Tao va dung network rieng.\strut
\end{minipage} & \begin{minipage}[t]{\linewidth}\raggedright
Chay docker network create ten-network roi recreate container tren network do.\strut
\end{minipage} \\
\cline{2-4}
 & 5.31: Khong chia se user namespace cua host & \begin{minipage}[t]{\linewidth}\raggedright
Bo host user namespace.\strut
\end{minipage} & \begin{minipage}[t]{\linewidth}\raggedright
Bo userns host; can nhac userns-remap o daemon neu phu hop.\strut
\end{minipage} \\
\cline{2-4}
 & 5.32: Khong mount Docker socket vao container & \begin{minipage}[t]{\linewidth}\raggedright
Go Docker socket bind mount.\strut
\end{minipage} & \begin{minipage}[t]{\linewidth}\raggedright
Bo /var/run/docker.sock khoi volumes; dung API TLS hoac cong cu khong can daemon socket.\strut
\end{minipage} \\
\hline
\multirow{2}{=}{6} & Tránh Image sprawl & Xóa bỏ các images không sử dụng thủ công & Chạy lệnh docker image prune -a --force \\
\cline{2-4}
& Tránh container sprawl & Xóa bỏ các containers đã dừng thủ công & Chạy lệnh docker container prune --force \\
\hline
\end{longtable}
}
